\documentclass[10pt,a4paper]{amsart}
\usepackage[margin=19mm]{geometry}
\usepackage{amsmath,amssymb,mathtools}
\usepackage[scr=rsfs,cal=euler]{mathalfa}
\usepackage{xfrac}
\usepackage[unicode,hidelinks,bookmarks=false]{hyperref}
\newcommand{\ie}{i.e.\ }
\newcommand{\cf}{cf.\ }
\newcommand{\resp}{respectively\ }
\newcommand{\Subsection}{\subsection}
\providecommand{\nobreakdash}{\nobreak\hbox{-}}
\setlength{\emergencystretch}{2em}
\multlinegap0pt
\usepackage{bm}
\usepackage{bbm}
\usepackage{dsfont}
\usepackage{xspace}
\usepackage{cancel}
\usepackage{mathtools}
\usepackage{amsthm}
\makeatletter
\@ifundefined{theorem}{\newtheorem{theorem}{Theorem}}{}
\@ifundefined{assumption}{\newtheorem{assumption}[theorem]{Assumption}}{}
\makeatother
\title[From Simple to Complex: Curriculum-Guided Physics-Informed N]{From Simple to Complex: Curriculum-Guided Physics-Informed Neural Networks via Gaussian Mixture Models}
\author{Jianan Yang, Yiran Wang, Shuai Li, Fujun Cao, Xuefei Yan, Junmin Liu}
\date{Source: arXiv:2605.19263v1 (CC BY 4.0)}
\begin{document}
\maketitle

\noindent\emph{Source and licence.} From Simple to Complex: Curriculum-Guided Physics-Informed Neural Networks via Gaussian Mixture Models. Version arXiv:2605.19263v1 (CC BY 4.0), distributed under the Creative Commons Attribution 4.0 International licence, as recorded in the arXiv licence metadata for this exact version and shown on the abstract page https://arxiv.org/abs/2605.19263v1. Full rights notice: licenses/arxiv-2605.19263-CC-BY-4.0.txt.\par\smallskip

\noindent\textit{Editorial scope.} Counted unit: the Theorem whose source label is \texttt{thm:weight\_equiv} in the article (source0.tex lines 1321--1390, source label \texttt{thm:weight\_equiv}), delivered together with the article's own complete original proof of it (source0.tex lines 1392--1547, one \texttt{proof} environment of 156 lines). No mathematics is written, repaired or re-derived: statement and proof are the article's own text line for line. Required context retained uncounted: eq:gmm\_posterior (lines 250--256); eq:easy\_hard\_weights (lines 306--310); eq:curriculum\_weight (lines 314--319); eq:precision\_factor (lines 327--333); eq:effective\_precision (lines 340--343); eq:combined\_weight (lines 353--357); eq:weight\_normalize (lines 374--379); ass:bounded\_variances (lines 1237--1251). Every definition, display and label of those lines is retained unchanged. Omitted from this package and not redistributed: 1--249, 257--305, 311--313, 320--326, 334--339, 344--352, 358--373, 380--1236, 1252--1320, 1391--1391, 1548--2507. Nothing inside a retained excerpt is omitted or altered: every retained body is delivered exactly as the article's lines are written, so no omission receipt of any kind accompanies this package. Citations: the retained excerpts carry no citation command at all, so no bibliography entry of the article is transcribed and no reference is redistributed. Reference adaptation: none was needed and none was made; every reference inside the delivered unit resolves inside it, so no unresolved reference or citation is produced. Printed slips kept literal: no printed word, symbol, hypothesis or number of the counted statement or of its proof was altered. Layout and class: the article's own class is \texttt{article} with packages including microtype, graphicx, booktabs, bbm, amsmath, amssymb, mathtools, amsthm, algorithm, algorithmic, amsfonts, multirow, subcaption, hyperref; the wrapper is the lane's pinned \texttt{amsart} editorial wrapper at 10pt on A4, which restates the theorem-like environments the retained text needs and adds the article's own macro definitions that the retained text needs (none), with extra wrapper packages (bm, bbm, dsfont, xspace, cancel, mathtools). The delivered numbering is the unit's own. Assets: no retained span contains a figure environment, a graphics-inclusion command, a PDF-page inclusion, a table or a TikZ picture; no image file is copied into this package, the package declares no asset dependency and no image file is redistributed with it. Licence: version arXiv:2605.19263v1 of this article is distributed under the Creative Commons Attribution 4.0 International licence, as recorded in the arXiv licence metadata for this exact version and shown on the abstract page https://arxiv.org/abs/2605.19263v1. A full rights notice is delivered at licenses/arxiv-2605.19263-CC-BY-4.0.txt. Characters: every retained excerpt is pure ASCII, so the delivered engine, XeLaTeX with its default Latin Modern text font, reports no missing character.
	\begin{equation}
		\gamma_{im}
		= \frac{\pi_m\,\mathcal{N}(r_i\mid\mu_m,\sigma_m^2)}
		{\sum_{l=1}^{K}\pi_l\,
			\mathcal{N}(r_i\mid\mu_l,\sigma_l^2)}.
		\label{eq:gmm_posterior}
	\end{equation}

	\begin{equation}
		w_m^{\mathrm{easy}} = \exp(-\beta\tilde{d}_m),\;
		w_m^{\mathrm{hard}} = \exp(-\beta(1 - \tilde{d}_m)),
		\label{eq:easy_hard_weights}
	\end{equation}

	\begin{equation}
		w_m^{\mathrm{CL}}(\tau)
		= (1-\tau)\,w_m^{\mathrm{easy}}
		+ \tau\,w_m^{\mathrm{hard}}.
		\label{eq:curriculum_weight}
	\end{equation}

	\begin{equation}
		v_m
		= \frac{(\sigma_m^2+\epsilon)^{-1}}
		{\max_{1\leqslant l\leqslant K}(\sigma_l^2+\epsilon)^{-1}}
		\;\in\;(0,1].
		\label{eq:precision_factor}
	\end{equation}

	\begin{equation}
		\tilde{v}_m(\tau) = (1-\tau)\,v_m + \tau.
		\label{eq:effective_precision}
	\end{equation}

	\begin{equation}
		w_m^{\mathrm{comp}}(\tau)
		= w_m^{\mathrm{CL}}(\tau)\cdot\tilde{v}_m(\tau).
		\label{eq:combined_weight}
	\end{equation}

	\begin{equation}
		w_i \leftarrow
		\frac{N_\Omega\cdot w_i}
		{\sum_{j=1}^{N_\Omega}w_j + \epsilon}.
		\label{eq:weight_normalize}
	\end{equation}

\begin{assumption}[Bounded GMM variances]
	\label{ass:bounded_variances}
	There exist constants
	$0<\underline{\sigma}^{2}\leqslant \overline{\sigma}^{2}<\infty$
	such that, at every GMM refresh step and for every Gaussian component
	$m$,
	\begin{equation}
		\underline{\sigma}^{2}
		\leqslant
		\sigma_m^{2}
		\leqslant
		\overline{\sigma}^{2}.
		\label{eq:app_variance_bound}
	\end{equation}
\end{assumption}

\begin{theorem}[Uniform equivalence between weighted and standard PDE losses]
	\label{thm:weight_equiv}
	Under Assumption~\ref{ass:bounded_variances}, define
	\begin{equation}
		\underline{v}
		:=
		\frac{\underline{\sigma}^{2}+\epsilon}
		{\overline{\sigma}^{2}+\epsilon}
		\in(0,1].
		\label{eq:app_v_under}
	\end{equation}
	Then, for every Gaussian component $m$ and every curriculum parameter
	$\tau\in[0,1]$, the component-level weight
	\begin{equation}
		w_m^{\mathrm{comp}}(\tau)
		=
		w_m^{\mathrm{CL}}(\tau)\,\tilde v_m(\tau)
		\label{eq:app_component_weight}
	\end{equation}
	satisfies
	\begin{equation}
		e^{-\beta}\underline{v}
		\leqslant
		w_m^{\mathrm{comp}}(\tau)
		\leqslant
		1.
		\label{eq:app_component_weight_bound}
	\end{equation}
	Define the pre-normalized sample weight by
	\begin{equation}
		w_i^{\mathrm{raw}}(\tau)
		:=
		\sum_{m=1}^{K}\gamma_{im}\,w_m^{\mathrm{comp}}(\tau).
		\label{eq:app_raw_weight}
	\end{equation}
	Then
	\begin{equation}
		e^{-\beta}\underline{v}
		\leqslant
		w_i^{\mathrm{raw}}(\tau)
		\leqslant
		1.
		\label{eq:app_raw_weight_bound}
	\end{equation}
	Let
	\begin{equation}
		\begin{aligned}
			c_- &:= \frac{N_{\Omega}e^{-\beta}\underline{v}}{N_{\Omega}+\epsilon},\\
			c_+ &:= \frac{N_{\Omega}}{N_{\Omega}e^{-\beta}\underline{v}+\epsilon}.
		\end{aligned}
		\label{eq:app_weight_constants}
	\end{equation}
	After the normalization in Eq.~\eqref{eq:weight_normalize}, the final
	sample weight $w_i$ satisfies
	\begin{equation}
		c_- \leqslant w_i \leqslant c_+.
		\label{eq:app_final_weight_bound}
	\end{equation}
	Consequently, for every $\theta$,
	\begin{equation}
		c_-\,
		\mathcal{L}_{\mathrm{PDE}}(\theta)
		\leqslant
		\mathcal{L}_{\mathrm{PDE}}^{w}(\theta)
		\leqslant
		c_+\,
		\mathcal{L}_{\mathrm{PDE}}(\theta).
		\label{eq:app_loss_equivalence}
	\end{equation}
\end{theorem}

\begin{proof}
	By Eq.~\eqref{eq:easy_hard_weights},
	\[
	w_m^{\mathrm{easy}}=\exp(-\beta\tilde d_m),
	\quad
	w_m^{\mathrm{hard}}=\exp(-\beta(1-\tilde d_m)).
	\]
	Since $\tilde d_m\in[0,1]$, one has
	\begin{equation}
		e^{-\beta}\leqslant w_m^{\mathrm{easy}}\leqslant 1,
		\quad
		e^{-\beta}\leqslant w_m^{\mathrm{hard}}\leqslant 1.
		\label{eq:app_easy_hard_bound}
	\end{equation}
	By Eq.~\eqref{eq:curriculum_weight},
	\[
	w_m^{\mathrm{CL}}(\tau)
	=
	(1-\tau)w_m^{\mathrm{easy}}
	+
	\tau w_m^{\mathrm{hard}},
	\]
	which is a convex combination of
	$w_m^{\mathrm{easy}}$ and $w_m^{\mathrm{hard}}$.
	Hence,
	\begin{equation}
		e^{-\beta}\leqslant w_m^{\mathrm{CL}}(\tau)\leqslant 1.
		\label{eq:app_cl_bound}
	\end{equation}
	
	Next, by Assumption~\ref{ass:bounded_variances},
	\begin{equation}
		\frac{1}{\overline{\sigma}^{2}+\epsilon}
		\leqslant
		(\sigma_m^{2}+\epsilon)^{-1}
		\leqslant
		\frac{1}{\underline{\sigma}^{2}+\epsilon}.
		\label{eq:app_inverse_variance_bound}
	\end{equation}
	Using Eq.~\eqref{eq:precision_factor}, we have
	\begin{align}
		v_m
		&=
		\frac{(\sigma_m^{2}+\epsilon)^{-1}}
		{\max_{1\leqslant l\leqslant K}(\sigma_l^{2}+\epsilon)^{-1}}
		\geqslant
		\frac{(\overline{\sigma}^{2}+\epsilon)^{-1}}
		{(\underline{\sigma}^{2}+\epsilon)^{-1}}
		=
		\frac{\underline{\sigma}^{2}+\epsilon}
		{\overline{\sigma}^{2}+\epsilon}
		=
		\underline{v}.
		\label{eq:app_vm_lower}
	\end{align}
	Trivially, $v_m\leqslant 1$. Then by Eq.~\eqref{eq:effective_precision},
	\begin{equation}
		\underline{v}\leqslant \tilde v_m(\tau)\leqslant 1.
		\label{eq:app_vtilde_bound}
	\end{equation}
	Combining \eqref{eq:app_cl_bound},
	\eqref{eq:app_vtilde_bound}, and
	Eq.~\eqref{eq:combined_weight}, we obtain
	\[
	e^{-\beta}\underline{v}
	\leqslant
	w_m^{\mathrm{comp}}(\tau)
	=
	w_m^{\mathrm{CL}}(\tau)\tilde v_m(\tau)
	\leqslant
	1,
	\]
	which proves \eqref{eq:app_component_weight_bound}.
	
	Now consider the pre-normalized sample weight
	\[
	w_i^{\mathrm{raw}}(\tau)
	=
	\sum_{m=1}^{K}\gamma_{im}\,w_m^{\mathrm{comp}}(\tau).
	\]
	Since $\gamma_{im}\geqslant 0$ and
	$\sum_{m=1}^{K}\gamma_{im}=1$ by Eq.~\eqref{eq:gmm_posterior},
	$w_i^{\mathrm{raw}}(\tau)$ is a convex combination of
	$\{w_m^{\mathrm{comp}}(\tau)\}_{m=1}^{K}$.
	Therefore,
	\[
	e^{-\beta}\underline{v}
	\leqslant
	w_i^{\mathrm{raw}}(\tau)
	\leqslant
	1,
	\]
	which proves \eqref{eq:app_raw_weight_bound}.
	
	By Eq.~\eqref{eq:weight_normalize}, the final sample weight is
	\begin{equation}
		w_i
		=
		\frac{N_{\Omega}\,w_i^{\mathrm{raw}}(\tau)}
		{\sum_{j=1}^{N_{\Omega}}w_j^{\mathrm{raw}}(\tau)+\epsilon}.
		\label{eq:app_exact_normalization}
	\end{equation}
	Using \eqref{eq:app_raw_weight_bound}, we have
	\[
	\sum_{j=1}^{N_{\Omega}}w_j^{\mathrm{raw}}(\tau)
	\leqslant
	N_{\Omega},
	\quad
	\sum_{j=1}^{N_{\Omega}}w_j^{\mathrm{raw}}(\tau)
	\geqslant
	N_{\Omega}e^{-\beta}\underline{v}.
	\]
	Substituting these bounds into \eqref{eq:app_exact_normalization} yields
	\[
	w_i
	\geqslant
	\frac{N_{\Omega}e^{-\beta}\underline{v}}{N_{\Omega}+\epsilon}
	=
	c_-,
	\quad
	w_i
	\leqslant
	\frac{N_{\Omega}}{N_{\Omega}e^{-\beta}\underline{v}+\epsilon}
	=
	c_+,
	\]
	which proves \eqref{eq:app_final_weight_bound}.
	
	Finally, since
	$r_{\mathrm{PDE},i}^{2}\geqslant 0$ for all $i$,
	\begin{align}
		\mathcal{L}_{\mathrm{PDE}}^{w}(\theta)
		&=
		\frac{1}{N_{\Omega}}
		\sum_{i=1}^{N_{\Omega}}
		w_i\,r_{\mathrm{PDE},i}^{2}
		\nonumber\\
		&\geqslant
		\frac{1}{N_{\Omega}}
		\sum_{i=1}^{N_{\Omega}}
		c_-\,r_{\mathrm{PDE},i}^{2}
		\nonumber\\
		&=
		c_-\mathcal{L}_{\mathrm{PDE}}(\theta),
		\label{eq:app_loss_lower}
	\end{align}
	and similarly,
	\begin{align}
		\mathcal{L}_{\mathrm{PDE}}^{w}(\theta)
		\leqslant
		c_+\mathcal{L}_{\mathrm{PDE}}(\theta).
		\label{eq:app_loss_upper}
	\end{align}
	Combining \eqref{eq:app_loss_lower} and \eqref{eq:app_loss_upper}
	proves \eqref{eq:app_loss_equivalence}.
\end{proof}

\end{document}
